\subsection{نتایج تجربی}
ارزیابی آزمایشگاهی ماژول \lr{admm\_compressor} بر روی چهار مدل مبنا نتایج قابل توجه مندرج در جدول~\ref{tab:benchmark_admm_compressor} را نشان می‌دهد:

\begin{table}[H]
\centering
\small
\caption{کارنامه تجربی و ارزیابی بنچ‌مارک سخت‌افزاری ماژول \lr{admm\_compressor} روی مدل‌های پایه}
\label{tab:benchmark_admm_compressor}
\begin{tabular}{lccccc}
\toprule
\textbf{مدل پایه \lr{(Base Model)}} & \textbf{دقت \lr{(Accuracy)}} & \textbf{شاخص \lr{F1}} & \textbf{حافظه \lr{(Peak VRAM)}} & \textbf{توان \lr{(Throughput)}} & \textbf{وضعیت} \\
\midrule
\lr{DistilBERT} & \lr{88.53\%} & \lr{0.8922} & \lr{339.42 MB} & \lr{54,607.47 tok/s} & \lr{PASS} \\
\lr{RoBERTa} & \lr{89.11\%} & \lr{0.8966} & \lr{333.44 MB} & \lr{13,290.18 tok/s} & \lr{PASS} \\
\lr{GPT-2} & \lr{84.63\%} & \lr{0.8537} & \lr{409.87 MB} & \lr{15,433.04 tok/s} & \lr{PASS} \\
\lr{Qwen} & \lr{88.42\%} & \lr{0.8947} & \lr{1,851.31 MB} & \lr{1,099.69 tok/s} & \lr{PASS} \\
\bottomrule
\end{tabular}
\end{table}

\subsection{محدودیت‌ها و تحلیل شکست}
تحلیل چندوجهی کارایی این ماژول نکات برجسته‌ای را آشکار می‌سازد:
\begin{itemize}
\item \textbf{پایداری همه‌جانبه با نمره قبولی کامل (PASS):} تمامی چهار مدل تحت آزمون موفق به کسب نرخ بالای دقت (بین \lr{84.63\%} تا \lr{89.11\%}) با همگرایی پایدار لاگرانژی شدند.
\item \textbf{تعادل قابل توجه دقت و مصرف حافظه:} مدل \lr{RoBERTa} با دقت \lr{89.11\%} و F1 معادل \lr{0.8966} و مصرف حافظه تنها \lr{333.44 MB}، گواهی بر قدرت همگرایی هم‌زمان هرس و کوانتیزاسیون است.
\item \textbf{چالش زمان کالیبراسیون و بار محاسباتی بهینه‌سازی:} علی‌رغم نتایج استنتاجی عالی، اجرای تکرارهای متناوب حل بسته و تصاویر غیرخطی نیازمند صرف زمان پردازشی اولیه و حافظه محاسباتی برای نگهداری متغیرهای دوگان ($\mathbf{U}_1, \mathbf{U}_2$) در طول فشرده‌سازی است. در صورتی که تعداد تکرارها کمتر از ۴ تنظیم شود، واگرایی دوگان رخ داده و خطای بازسازی به شدت افزایش می‌یابد.
\end{itemize}
